\section{Algorithm}
\label{sec:algorithm}

This section specifies the agentic system built on \textsc{DeXposure-FM} for two tasks: network risk monitoring (Section~\ref{subsec:task_monitoring}--\ref{subsec:scenario_engine}) and decision recommendation (Section~\ref{subsec:decision_v1}).

Figure~\ref{fig:system_overview} provides an architectural overview of the full system.

\begin{figure*}[!t]
    \centering
    \resizebox{0.95\textwidth}{!}{% Fig. 1 -- DeXposure-Claw block diagram (TikZ, TKDE classic style)
% Drop-in replacement for figures/fig1_system_overview.{pdf,png}.
% Loaded from sections/Alg.tex via % Fig. 1 -- DeXposure-Claw block diagram (TikZ, TKDE classic style)
% Drop-in replacement for figures/fig1_system_overview.{pdf,png}.
% Loaded from sections/Alg.tex via % Fig. 1 -- DeXposure-Claw block diagram (TikZ, TKDE classic style)
% Drop-in replacement for figures/fig1_system_overview.{pdf,png}.
% Loaded from sections/Alg.tex via \input{figures/fig1_system_overview}
%
% Required packages (already in the main paper tex):
%   \usepackage{tikz}
%
\begin{tikzpicture}[
  font=\small,
  layer/.style    = {draw, thick, rectangle, minimum width=2.8cm, minimum height=2.6cm,
                     align=center, inner sep=4pt},
  io/.style       = {draw, thick, rectangle, minimum width=2.0cm, minimum height=2.6cm,
                     align=center, inner sep=4pt},
  flow/.style     = {-{Latex[length=2.2mm,width=1.8mm]}, thick},
  tag/.style      = {font=\scriptsize\sffamily, text=black!55},
  hint/.style     = {font=\scriptsize\itshape, text=black!55},
  rule/.style     = {draw=black!22, line width=0.35pt},
  ablation/.style = {font=\scriptsize\sffamily\bfseries, text=red!75!black},
]

% --- input box -------------------------------------------------------------
\node[io] (in) {%
  \scriptsize\sffamily INPUT \\[2pt]
  \rule{1.4cm}{0.3pt} \\[6pt]
  $G_t$ \\[2pt]
  \scriptsize\textit{weekly exposure} \\
  \scriptsize\textit{graph snapshot}
};

% --- four layers -----------------------------------------------------------
\node[layer, right=0.55cm of in]   (l1) {%
  {\scriptsize\sffamily LAYER 1} \\[2pt]
  {\textbf{FM Prediction}} \\[3pt]
  {\scriptsize\itshape DeXposure-FM} \\
  {\scriptsize\itshape (GraphPFN hybrid)} \\[6pt]
  \rule{2.2cm}{0.3pt} \\[3pt]
  {\scriptsize OUT} \\
  {\small $\hat G_{t+h},\; \tilde G^{(1..S)}_{t+h}$}
};

\node[layer, right=0.55cm of l1]   (l2) {%
  {\scriptsize\sffamily LAYER 2} \\[2pt]
  {\textbf{Monitor \& Scenario}} \\[3pt]
  {\scriptsize\itshape 5 metrics + z-scores} \\
  {\scriptsize\itshape 5 scenarios $S_1..S_5$} \\[6pt]
  \rule{2.2cm}{0.3pt} \\[3pt]
  {\scriptsize OUT} \\
  {\small alerts $\mathcal{A}_t$,\; losses $\mathcal{L}_t$}
};

\node[layer, right=0.55cm of l2]   (l3) {%
  {\scriptsize\sffamily LAYER 3} \\[2pt]
  {\textbf{LLM Decision}} \\[3pt]
  {\scriptsize\itshape Claude Opus 4.7} \\
  {\scriptsize\itshape reasons over evidence} \\[6pt]
  \rule{2.2cm}{0.3pt} \\[3pt]
  {\scriptsize OUT} \\
  {\small draft ticket $\tau$}
};

\node[layer, right=0.55cm of l3]   (l4) {%
  {\scriptsize\sffamily LAYER 4} \\[2pt]
  {\textbf{Safety Gate}} \\[3pt]
  {\scriptsize\itshape DH$_t \geq \tau_{\mathrm{data}}$} \\
  {\scriptsize\itshape C$_t \geq \tau_{\mathrm{conf}}$} \\[6pt]
  \rule{2.2cm}{0.3pt} \\[3pt]
  {\scriptsize OUT} \\
  {\small $\tau^\star$ or SafeMode}
};

% --- output ---------------------------------------------------------------
\node[io, right=0.55cm of l4] (out) {%
  \scriptsize\sffamily OUTPUT \\[2pt]
  \rule{1.4cm}{0.3pt} \\[6pt]
  $\tau^\star$ \\[2pt]
  \scriptsize\textit{supervisory} \\
  \scriptsize\textit{ticket}
};

% --- arrows --------------------------------------------------------------
\draw[flow] (in.east)  -- (l1.west);
\draw[flow] (l1.east)  -- (l2.west);
\draw[flow] (l2.east)  -- (l3.west);
\draw[flow] (l3.east)  -- (l4.west);
\draw[flow] (l4.east)  -- (out.west);

% --- bottom annotation row: layer-wise ablation lifts -------------------
\node[hint, below=4pt of l1, anchor=north] {\textit{vs persistence: trend $\uparrow$}};
\node[hint, below=4pt of l2, anchor=north] {\textit{CVaR aggregation}};
\node[ablation, below=4pt of l3, anchor=north] {F1 +27\% over rules};
\node[ablation, below=4pt of l4, anchor=north] {judge +0.21};

\end{tikzpicture}

%
% Required packages (already in the main paper tex):
%   \usepackage{tikz}
%
\begin{tikzpicture}[
  font=\small,
  layer/.style    = {draw, thick, rectangle, minimum width=2.8cm, minimum height=2.6cm,
                     align=center, inner sep=4pt},
  io/.style       = {draw, thick, rectangle, minimum width=2.0cm, minimum height=2.6cm,
                     align=center, inner sep=4pt},
  flow/.style     = {-{Latex[length=2.2mm,width=1.8mm]}, thick},
  tag/.style      = {font=\scriptsize\sffamily, text=black!55},
  hint/.style     = {font=\scriptsize\itshape, text=black!55},
  rule/.style     = {draw=black!22, line width=0.35pt},
  ablation/.style = {font=\scriptsize\sffamily\bfseries, text=red!75!black},
]

% --- input box -------------------------------------------------------------
\node[io] (in) {%
  \scriptsize\sffamily INPUT \\[2pt]
  \rule{1.4cm}{0.3pt} \\[6pt]
  $G_t$ \\[2pt]
  \scriptsize\textit{weekly exposure} \\
  \scriptsize\textit{graph snapshot}
};

% --- four layers -----------------------------------------------------------
\node[layer, right=0.55cm of in]   (l1) {%
  {\scriptsize\sffamily LAYER 1} \\[2pt]
  {\textbf{FM Prediction}} \\[3pt]
  {\scriptsize\itshape DeXposure-FM} \\
  {\scriptsize\itshape (GraphPFN hybrid)} \\[6pt]
  \rule{2.2cm}{0.3pt} \\[3pt]
  {\scriptsize OUT} \\
  {\small $\hat G_{t+h},\; \tilde G^{(1..S)}_{t+h}$}
};

\node[layer, right=0.55cm of l1]   (l2) {%
  {\scriptsize\sffamily LAYER 2} \\[2pt]
  {\textbf{Monitor \& Scenario}} \\[3pt]
  {\scriptsize\itshape 5 metrics + z-scores} \\
  {\scriptsize\itshape 5 scenarios $S_1..S_5$} \\[6pt]
  \rule{2.2cm}{0.3pt} \\[3pt]
  {\scriptsize OUT} \\
  {\small alerts $\mathcal{A}_t$,\; losses $\mathcal{L}_t$}
};

\node[layer, right=0.55cm of l2]   (l3) {%
  {\scriptsize\sffamily LAYER 3} \\[2pt]
  {\textbf{LLM Decision}} \\[3pt]
  {\scriptsize\itshape Claude Opus 4.7} \\
  {\scriptsize\itshape reasons over evidence} \\[6pt]
  \rule{2.2cm}{0.3pt} \\[3pt]
  {\scriptsize OUT} \\
  {\small draft ticket $\tau$}
};

\node[layer, right=0.55cm of l3]   (l4) {%
  {\scriptsize\sffamily LAYER 4} \\[2pt]
  {\textbf{Safety Gate}} \\[3pt]
  {\scriptsize\itshape DH$_t \geq \tau_{\mathrm{data}}$} \\
  {\scriptsize\itshape C$_t \geq \tau_{\mathrm{conf}}$} \\[6pt]
  \rule{2.2cm}{0.3pt} \\[3pt]
  {\scriptsize OUT} \\
  {\small $\tau^\star$ or SafeMode}
};

% --- output ---------------------------------------------------------------
\node[io, right=0.55cm of l4] (out) {%
  \scriptsize\sffamily OUTPUT \\[2pt]
  \rule{1.4cm}{0.3pt} \\[6pt]
  $\tau^\star$ \\[2pt]
  \scriptsize\textit{supervisory} \\
  \scriptsize\textit{ticket}
};

% --- arrows --------------------------------------------------------------
\draw[flow] (in.east)  -- (l1.west);
\draw[flow] (l1.east)  -- (l2.west);
\draw[flow] (l2.east)  -- (l3.west);
\draw[flow] (l3.east)  -- (l4.west);
\draw[flow] (l4.east)  -- (out.west);

% --- bottom annotation row: layer-wise ablation lifts -------------------
\node[hint, below=4pt of l1, anchor=north] {\textit{vs persistence: trend $\uparrow$}};
\node[hint, below=4pt of l2, anchor=north] {\textit{CVaR aggregation}};
\node[ablation, below=4pt of l3, anchor=north] {F1 +27\% over rules};
\node[ablation, below=4pt of l4, anchor=north] {judge +0.21};

\end{tikzpicture}

%
% Required packages (already in the main paper tex):
%   \usepackage{tikz}
%
\begin{tikzpicture}[
  font=\small,
  layer/.style    = {draw, thick, rectangle, minimum width=2.8cm, minimum height=2.6cm,
                     align=center, inner sep=4pt},
  io/.style       = {draw, thick, rectangle, minimum width=2.0cm, minimum height=2.6cm,
                     align=center, inner sep=4pt},
  flow/.style     = {-{Latex[length=2.2mm,width=1.8mm]}, thick},
  tag/.style      = {font=\scriptsize\sffamily, text=black!55},
  hint/.style     = {font=\scriptsize\itshape, text=black!55},
  rule/.style     = {draw=black!22, line width=0.35pt},
  ablation/.style = {font=\scriptsize\sffamily\bfseries, text=red!75!black},
]

% --- input box -------------------------------------------------------------
\node[io] (in) {%
  \scriptsize\sffamily INPUT \\[2pt]
  \rule{1.4cm}{0.3pt} \\[6pt]
  $G_t$ \\[2pt]
  \scriptsize\textit{weekly exposure} \\
  \scriptsize\textit{graph snapshot}
};

% --- four layers -----------------------------------------------------------
\node[layer, right=0.55cm of in]   (l1) {%
  {\scriptsize\sffamily LAYER 1} \\[2pt]
  {\textbf{FM Prediction}} \\[3pt]
  {\scriptsize\itshape DeXposure-FM} \\
  {\scriptsize\itshape (GraphPFN hybrid)} \\[6pt]
  \rule{2.2cm}{0.3pt} \\[3pt]
  {\scriptsize OUT} \\
  {\small $\hat G_{t+h},\; \tilde G^{(1..S)}_{t+h}$}
};

\node[layer, right=0.55cm of l1]   (l2) {%
  {\scriptsize\sffamily LAYER 2} \\[2pt]
  {\textbf{Monitor \& Scenario}} \\[3pt]
  {\scriptsize\itshape 5 metrics + z-scores} \\
  {\scriptsize\itshape 5 scenarios $S_1..S_5$} \\[6pt]
  \rule{2.2cm}{0.3pt} \\[3pt]
  {\scriptsize OUT} \\
  {\small alerts $\mathcal{A}_t$,\; losses $\mathcal{L}_t$}
};

\node[layer, right=0.55cm of l2]   (l3) {%
  {\scriptsize\sffamily LAYER 3} \\[2pt]
  {\textbf{LLM Decision}} \\[3pt]
  {\scriptsize\itshape Claude Opus 4.7} \\
  {\scriptsize\itshape reasons over evidence} \\[6pt]
  \rule{2.2cm}{0.3pt} \\[3pt]
  {\scriptsize OUT} \\
  {\small draft ticket $\tau$}
};

\node[layer, right=0.55cm of l3]   (l4) {%
  {\scriptsize\sffamily LAYER 4} \\[2pt]
  {\textbf{Safety Gate}} \\[3pt]
  {\scriptsize\itshape DH$_t \geq \tau_{\mathrm{data}}$} \\
  {\scriptsize\itshape C$_t \geq \tau_{\mathrm{conf}}$} \\[6pt]
  \rule{2.2cm}{0.3pt} \\[3pt]
  {\scriptsize OUT} \\
  {\small $\tau^\star$ or SafeMode}
};

% --- output ---------------------------------------------------------------
\node[io, right=0.55cm of l4] (out) {%
  \scriptsize\sffamily OUTPUT \\[2pt]
  \rule{1.4cm}{0.3pt} \\[6pt]
  $\tau^\star$ \\[2pt]
  \scriptsize\textit{supervisory} \\
  \scriptsize\textit{ticket}
};

% --- arrows --------------------------------------------------------------
\draw[flow] (in.east)  -- (l1.west);
\draw[flow] (l1.east)  -- (l2.west);
\draw[flow] (l2.east)  -- (l3.west);
\draw[flow] (l3.east)  -- (l4.west);
\draw[flow] (l4.east)  -- (out.west);

% --- bottom annotation row: layer-wise ablation lifts -------------------
\node[hint, below=4pt of l1, anchor=north] {\textit{vs persistence: trend $\uparrow$}};
\node[hint, below=4pt of l2, anchor=north] {\textit{CVaR aggregation}};
\node[ablation, below=4pt of l3, anchor=north] {F1 +27\% over rules};
\node[ablation, below=4pt of l4, anchor=north] {judge +0.21};

\end{tikzpicture}
}
    \caption{\textsc{DeXposure-Claw}: a four-layer query architecture. The system ingests the weekly exposure graph $G_t$, runs FM prediction (Layer~1), monitoring and scenario engines (Layer~2), an LLM decision layer (Layer~3), and a safety gate (Layer~4) that emits either the supervisory ticket $\tau^\star$ or a safe-mode notice.}
    \label{fig:system_overview}
\end{figure*}

\subsection{System Interface}
\label{subsec:system_interface}
At each decision epoch $t$, the agent observes a credit-exposure snapshot $G_t=(V,E_t,W_t)$ and node covariates $X_t$, and interacts with \textsc{DeXposure-FM} through a single forecast primitive:
\begin{equation}
\label{eq:agent_forecast_primitive}
\mathcal{P}_{t,h} \leftarrow \textsc{DeXposure-FM}.\mathrm{Forecast}(G_t,X_t,h), \quad \forall h\in\mathcal{H},
\end{equation}
where $\mathcal{P}_{t,h}$ denotes the model's predictive distribution over the future network $G_{t+h}$, and $\mathcal{H}=\{1,4,8,12\}$ (weeks). From $\mathcal{P}_{t,h}$ we construct a representative predicted graph $\widehat{G}_{t+h}$ using a thresholded expected-weight operator:
\begin{align}
\label{eq:expected_graph}
\widehat{w}_{pq,t+h}
&:=
\Pr(e_{pq,t+h}=1)\cdot \mathbb{E}[w_{pq,t+h}\mid e_{pq,t+h}=1],
\\
\widehat{E}_{t+h}
&:=\{(p,q): \Pr(e_{pq,t+h}=1)>\pi_{\min}\}.
\end{align}
where $\pi_{\min}$ (default $0.2$, tuned on validation) balances edge coverage and sparsity. For uncertainty quantification, we draw $M$ Monte Carlo samples $\{\widetilde{G}^{(m)}_{t+h}\}_{m=1}^M$ by applying Bernoulli edge sampling and Gaussian weight perturbation to $\mathcal{P}_{t,h}$ (default $M=50$).

The agent returns three outputs: (a) a set of monitoring alerts with evidence and confidence, (b) a scenario stress summary, and (c) a ranked list of recommendation tickets.

\subsection{Data-Health Gating}
\label{subsec:data_health_gate}
Real-world DeFi graphs can be sparse, stale, or corrupted. Before producing actionable outputs, the agent computes a scalar data-health score from four deterministic checks---freshness, feature missingness, topology sanity (edge-to-node ratio), and discontinuity detection:
\begin{align}
\label{eq:data_health}
&DH_t := \mathrm{mean}\bigl(\mathrm{fresh}_t,\;\mathrm{miss}_t,\;\mathrm{topo}_t,\;\mathrm{disc}_t\bigr)\in[0,1],\nonumber\\
&\mathrm{SAFE\_MODE}_t := \mathbb{I}\{DH_t<\tau_{\mathrm{data}}\}.
\end{align}
When $DH_t$ falls below $\tau_{\mathrm{data}}$ (default $0.7$), the system enters safe mode: monitoring alerts are still emitted, but intervention-type recommendations are suppressed (Section~\ref{subsec:decision_v1}).

\subsection{Task~I: Network Risk Monitoring}
\label{subsec:task_monitoring}

\paragraph{Monitoring metrics ($\Phi$)}
Risk in exposure networks manifests as concentration (a few protocols dominate), fragility (shocks propagate widely), and instability (structure shifts rapidly). To capture these dimensions, the monitoring agent evaluates five functionals on $\widehat{G}_{t+h}$:

\begin{table*}[!t]
\centering
\footnotesize
\setlength{\tabcolsep}{4pt}
\begin{tabularx}{\textwidth}{@{}llX@{}}
\toprule
ID & Metric & Definition \\
\midrule
N1 & Systemic Importance & $\max_{p\in V} \mathrm{PageRank}(\widehat{G}_{t+h})_p$ \\
N2 & HHI Concentration & $\sum_p \bigl(d_p / \textstyle\sum_q d_q\bigr)^2$, \ $d_p = \text{weighted out-degree}$ \\
N3 & Network Density & $|E_{t+h}| \;/\; |V|(|V|-1)$ \\
N4 & PageRank Gini & $\mathrm{Gini}\!\bigl(\{\mathrm{PR}_p\}_{p\in V}\bigr)$ \\
N5 & Degree Gini & $\mathrm{Gini}\!\bigl(\{d_p\}_{p\in V}\bigr)$ \\
\bottomrule
\end{tabularx}
\caption{Monitoring functionals $\Phi$. PageRank uses power iteration with damping $0.85$ and $10$ iterations. Gini coefficient: $1 - 2\sum_{i=1}^{n}\sum_{k=1}^{i}x_{(k)} \;/\; (n \cdot \sum_k x_{(k)})$, where $x_{(k)}$ are sorted values.}
\label{tab:metrics}
\end{table*}

\noindent Formally, at each horizon $h$ we compute:
\begin{equation}
\label{eq:monitor_metrics}
\mathbf{N}_{t,h} := \Phi(\widehat{G}_{t+h},X_t)\in\mathbb{R}^5,\qquad
\widetilde{\mathbf{N}}^{(m)}_{t,h} := \Phi(\widetilde{G}^{(m)}_{t+h},X_t).
\end{equation}

\paragraph{Alert generation}
An alert fires when a predicted metric departs from its rolling historical baseline by more than $z$ standard deviations (default $z=2$, baseline window $L=26$ weeks):
\begin{equation}
\label{eq:alert_trigger}
\mathcal{A}_t :=
\Bigl\{\, a=(h,j)\;:\; \frac{|N_{t,h}^{(j)} - \mu_t^{(j)}|}{\max(\sigma_t^{(j)},\,\epsilon)} > z \Bigr\},
\end{equation}
where $(\mu_t^{(j)},\sigma_t^{(j)})$ are the rolling mean and standard deviation of metric $j$ over the last $L$ observed snapshots, and $\epsilon=10^{-8}$ prevents division by zero.

\paragraph{Evidence and attribution}
Each alert $a\in\mathcal{A}_t$ carries an evidence bundle $\mathrm{Ev}_t(a) = ( N_{t,h}^{(j)},\, \mu_t^{(j)},\, \sigma_t^{(j)},\, \mathrm{Attr}_t(a),\, h )$, where $\mathrm{Attr}_t(a)$ is a ranked list of top-$K$ contributing nodes/edges obtained via marginal contribution approximation: $\mathrm{Contrib}(e) := \Phi^{(j)}(\widehat{G}_{t+h}) - \Phi^{(j)}(\widehat{G}_{t+h}\setminus e)$ (default $K=10$).

\paragraph{Uncertainty-aware confidence}
Alert reliability depends on three factors: data quality, forecast dispersion, and horizon length. Using Monte Carlo samples, we compute the coefficient of variation $\mathrm{CV}_{t,h}^{(j)} := \sigma(\{\widetilde{N}^{(m,j)}_{t,h}\}_m) / |\bar{N}^{(j)}_{t,h}|$ (clamped to $[0,1]$) and define:
\begin{equation}
\label{eq:alert_confidence}
C_t(a)\;:=\; DH_t \cdot \frac{1}{1+\mathrm{CV}_{t,h}^{(j)}} \cdot \frac{1}{1+\ln(1+h)}\;\in[0,1].
\end{equation}
This score monotonically decreases with lower data health, wider predictive spread, and longer horizons, ensuring that unreliable alerts are automatically downweighted.

\subsection{Scenario Engine}
\label{subsec:scenario_engine}
To quantify tail risk beyond metric-level alerts, the agent evaluates a standardised stress-scenario library $\mathcal{S}_0$. Each scenario applies a deterministic shock---removing or reducing edge weights of targeted nodes---and propagates losses through the network:

\begin{table}[!t]
\centering
\footnotesize
\setlength{\tabcolsep}{4pt}
\begin{tabular}{llll}
\toprule
ID & Scenario & Target & Shock \\
\midrule
S1 & Single protocol failure & Highest-weight node & 100\% \\
S2 & Bridge cluster failure & All bridge-category nodes & 100\% \\
S3 & Stablecoin de-peg & All stablecoin nodes & 50\% \\
S4 & Sector-wide shock & All lending nodes & 30\% \\
S5 & Correlated stress & Top-10 nodes by weight & 20\% \\
\bottomrule
\end{tabular}
\caption{Standardised stress-scenario library $\mathcal{S}_0$.}
\label{tab:scenarios}
\end{table}

\noindent For each scenario $s$ and horizon $h$, we compute the systemic loss on both the predicted graph and each MC sample:
\begin{align}
\label{eq:scenario_loss}
L_{t,h}(s)
&:= \frac{\sum_e w_e(\widehat{G}_{t+h}) - \sum_e w_e(\widehat{G}_{t+h}^s)}
{\sum_e w_e(\widehat{G}_{t+h})},
\\
\widetilde{L}^{(m)}_{t,h}(s)
&:= \ell\!\left(\widetilde{G}^{(m)}_{t+h};\, s\right).
\end{align}
where $\widehat{G}_{t+h}^s$ denotes the network after applying shock $s$. The scenario engine also tracks the number of \emph{distressed} nodes (those losing $>50\%$ of total edge weight) and a propagation depth estimate $\lfloor\log_2(\text{affected nodes}+1)\rfloor + 1$.

\paragraph{CVaR aggregation}
To capture tail risk, we compute an empirical CVaR from the MC sample losses. Let $\lambda\in[0,1]$ (default $0.5$) control the tail fraction. We sort the $M$ sample losses in descending order and average the worst $\lceil \lambda M \rceil$:
\begin{equation}
\label{eq:cvar}
\mathrm{CVaR}_\lambda(s,h) := \frac{1}{\lceil \lambda M \rceil} \sum_{i=1}^{\lceil \lambda M \rceil} \widetilde{L}^{(i)}_{t,h}(s),
\end{equation}
where $\widetilde{L}^{(1)} \ge \widetilde{L}^{(2)} \ge \cdots$ are order statistics. The scenario summary $\mathcal{R}_t$ ranks scenarios by $L_{t,h}(s)$ and retains the worst-case scenario across all horizons.

\subsection{Task~II: Decision Agent}
\label{subsec:decision_v1}

\paragraph{Action space}
The decision module maps alerts and scenario results to recommendation tickets using a conservative playbook with four action types, ordered by escalation severity:

\begin{table}[!t]
\centering
\footnotesize
\setlength{\tabcolsep}{4pt}
\begin{tabular}{lcccc}
\toprule
Action & Severity & Weight $w_u$ & Safe-mode OK? & Min.\ alerts \\
\midrule
\texttt{Monitor} & Low & 0.25 & Yes & 1 \\
\texttt{Investigate} & Medium & 0.50 & Yes & 1 \\
\makecell[l]{\texttt{Recommend}\\[-1pt]\texttt{-Reduce}} & High & 0.75 & No & 2 \\
\texttt{Contingency} & Critical & 1.00 & No & 3 \\
\bottomrule
\end{tabular}
\caption{Playbook action space $\mathcal{U}$ with severity weights and feasibility constraints.}
\label{tab:playbook}
\end{table}

\paragraph{Feasibility constraints}
An intervention-type action $u$ (those requiring safe-mode off) is feasible only when both the data-health gate and the confidence gate are satisfied:
\begin{equation}
\label{eq:decision_constraints}
u \in \mathcal{C}_t
\quad \Longleftrightarrow \quad
\bigl(\mathrm{SAFE\_MODE}_t=0\bigr)\ \wedge\
\Bigl(\bar{C}_t \ge \tau_{\mathrm{conf}}\Bigr),
\end{equation}
where $\bar{C}_t = \frac{1}{|\mathcal{A}_t|}\sum_{a\in\mathcal{A}_t} C_t(a)$ is the mean alert confidence and $\tau_{\mathrm{conf}}$ (default $0.6$) gates interventions. Non-interventional actions (\texttt{Monitor}, \texttt{Investigate}) are always feasible when at least one alert exists.

\paragraph{Ticket scoring}
Each feasible recommendation is scored by a multiplicative heuristic combining playbook severity, alert confidence, and scenario impact:
\begin{equation}
\label{eq:ticket_score}
\mathrm{Score}_t(u) :=
w_u \cdot \bar{C}_t \cdot \mathrm{Imp}_t,
\end{equation}
where $w_u$ is the severity weight (Table~\ref{tab:playbook}), $\bar{C}_t$ is the mean alert confidence, and $\mathrm{Imp}_t := \frac{1}{|\mathcal{S}_0|}\sum_{s\in\mathcal{S}_0} \mathrm{CVaR}_\lambda(s)$ aggregates scenario impact as the mean CVaR across all scenarios. When no scenarios are available, $\mathrm{Imp}_t$ defaults to $1.0$ (neutral multiplier), ensuring alert-only tickets receive non-zero scores. The agent returns the top-$N$ tickets (default $N=10$) ranked by $\mathrm{Score}_t(u)$.

\paragraph{Target selection}
For intervention tickets, targets are selected by collecting the top-$K$ attributed nodes from all triggering alerts and worst-case scenarios, ranked by frequency-weighted vote.

\subsection{End-to-End Procedure}
Algorithm~\ref{alg:agent_loop} summarises the full agent loop.

\begin{algorithm}[!t]
\caption{DeXposure-Claw: one epoch}
\label{alg:agent_loop}
\begin{algorithmic}[1]
\Require State $(G_t,X_t)$; horizons $\mathcal{H}$; scenarios $\mathcal{S}_0$; thresholds $(\tau_{\mathrm{data}},\tau_{\mathrm{conf}},z,\pi_{\min})$; MC count $M$
\Ensure Alerts $\mathcal{A}_t$; scenario summary $\mathcal{R}_t$; tickets $\mathcal{D}_t$
\State $DH_t \gets \mathrm{DataHealth}(G_t,X_t)$;\; $\mathrm{SAFE\_MODE}_t \gets \mathbb{I}\{DH_t<\tau_{\mathrm{data}}\}$
\ForAll{$h\in\mathcal{H}$}
    \State $\mathcal{P}_{t,h} \gets \textsc{DeXposure-FM}.\mathrm{Forecast}(G_t,X_t,h)$
    \State $\widehat{G}_{t+h} \gets \mathrm{BuildPredGraph}(\mathcal{P}_{t,h};\,\pi_{\min})$ \Comment{Eq.~\eqref{eq:expected_graph}}
    \State $\{\widetilde{G}^{(m)}_{t+h}\}_{m=1}^M \gets \mathrm{MCSample}(\mathcal{P}_{t,h})$ \Comment{Bernoulli edges + Gaussian weights}
    \State $\mathbf{N}_{t,h} \gets \Phi(\widehat{G}_{t+h},X_t)$;\; $\widetilde{\mathbf{N}}^{(m)}_{t,h} \gets \Phi(\widetilde{G}^{(m)}_{t+h},X_t)$ \Comment{Table~\ref{tab:metrics}: N1..N5}
    \State $(\mu^{(j)}_t,\sigma^{(j)}_t) \gets \mathrm{RollingStats}(\{\mathbf{N}_{t-l,h}\}_{l=1}^{L})$ \Comment{$L=26$ weeks}
    \State $\mathcal{A}_{t,h} \gets \{(h,j): |N^{(j)}_{t,h}-\mu^{(j)}_t|/\max(\sigma^{(j)}_t,\epsilon) > z\}$ \Comment{Eq.~\eqref{eq:alert_trigger}}
    \State Attach $\mathrm{Ev}_t(a)$ and $C_t(a)$ for all $a\in\mathcal{A}_{t,h}$ \Comment{Eq.~\eqref{eq:alert_confidence}, uses $\widetilde{\mathbf{N}}^{(m)}_{t,h}$ for CV}
    \ForAll{$s\in\mathcal{S}_0$}
        \State $L_{t,h}(s) \gets \mathrm{StressLoss}(\widehat{G}_{t+h}, s)$;\; $\widetilde{L}^{(m)}_{t,h}(s) \gets \mathrm{StressLoss}(\widetilde{G}^{(m)}_{t+h}, s)$
        \State $\mathrm{CVaR}_\lambda(s,h) \gets \mathrm{TailMean}(\{\widetilde{L}^{(m)}_{t,h}(s)\}_{m=1}^M, \lambda)$ \Comment{Eq.~\eqref{eq:cvar}}
    \EndFor
\EndFor
\State $\mathcal{A}_t \gets \bigcup_{h\in\mathcal{H}}\mathcal{A}_{t,h}$
\State $\mathcal{R}_t \gets \mathrm{RankScenarios}(\{L_{t,h}(s),\mathrm{CVaR}_\lambda(s,h)\}_{h,s})$
\State $\mathcal{D}_t \gets \mathrm{PlaybookDecision}($
\Statex \hspace{\algorithmicindent}$\mathcal{A}_t,\mathcal{R}_t,DH_t,\mathrm{SAFE\_MODE}_t,\tau_{\mathrm{conf}})$ \Comment{Table~\ref{tab:playbook}}
\State \Return $(\mathcal{A}_t,\mathcal{R}_t,\mathcal{D}_t)$
\end{algorithmic}
\end{algorithm}
